% Preamble template for Tufte-style lecture notes.
% Do not compile this file directly; compile the main notes file instead.
\documentclass{tufte-handout}

\usepackage{ntheorem}
\usepackage{graphicx}
\usepackage{amsmath}
\usepackage{amssymb}
\usepackage{hyperref}
\usepackage{epigraph}
\usepackage{booktabs}
\theoremstyle{break}
% \usepackage[
% bibencoding=utf8,% .bib file encoding
% maxbibnames=3, % otherwise et al
% minbibnames=1, % otherwise et al
% backend=biber,%
% sortlocale=en_US,%
% style=apa,% or authoryear
% % apabackref=false, % backreferences
% natbib=true,% for citet/citep, but this is for backward compatibility
% uniquename=false,%
% url=true,%
% sortcites=false,
% doi=true,%
% eprint=true%
% ]{biblatex}
% \addbibresource{lecture_note_bib.bib}

\hypersetup{
  colorlinks,
  urlcolor = blue,
  pdfauthor={Paul Goldsmith-Pinkham}
  pdfkeywords={econometrics}
  pdftitle={Lecture Notes for Applied Empirical Methods}
  pdfpagemode=UseNone
}
\newtheorem{ruleN}{Rule}
\newtheorem{thmN}{Theorem}
\newtheorem{assN}{Assumption}
\newtheorem{defN}{Definition}
\newtheorem{exmp}{Example}
\newtheorem{cmt}{Comment}
\newtheorem{proof}{Proof}

\newcommand{\continuation}{??}
\newtheorem*{excont}{Example \continuation}
\newenvironment{continueexample}[1]
 {\renewcommand{\continuation}{\ref{#1}}\excont[continued]}
 {\endexcont}
\newcommand{\bY}{\mathbf{Y}}
\newcommand{\bX}{\mathbf{X}}
\newcommand{\bD}{\mathbf{D}}
\newcommand{\E}{\mathbb{E}}

\newcommand\independent{\protect\mathpalette{\protect\independenT}{\perp}}
\def\independenT#1#2{\mathrel{\rlap{$#1#2$}\mkern2mu{#1#2}}}
\DeclareMathOperator{\Supp}{Supp}

\usepackage{cleveref}
\crefname{appsec}{appendix}{appendices}
\crefname{appsubsec}{appendix}{appendices}
\crefname{assumption}{assumption}{assumptions}
\crefname{equation}{equation}{equations}
\crefname{exmp}{example}{examples}
\crefname{assN}{assumption}{assumptions}
\crefname{cmt}{comment}{comments}
\crefname{defN}{definition}{definitions}

\usepackage[nolist]{acronym}
\begin{acronym}
  \acro{CI}{confidence interval}%
  \acro{OLS}{ordinary least squares}%
  \acro{CLT}{central limit theorem}%
  \acro{IV}{instrumental variables}%
  \acro{ATE}{average treatment effect}%
  \acro{RCT}{randomized control trial}%
  \acro{SUTVA}{stable unit treatment value assignment}
  \acro{VAM}{value-added model}%
  \acro{LAN}{locally asymptotically normal}%
  \acro{DiD}{difference-in-differences}%
  \acro{OVB}{omitted variables bias}
  \acro{FWL}{Frisch-Waugh-Lovell}
  \acro{DAG}{directed acyclic graph}
  \acro{PO}{potential outcomes}
\end{acronym}

\def\inprobHIGH{\,{\buildrel p \over \rightarrow}\,} 
\def\inprob{\,{\inprobHIGH}\,} 
\def\indistHIGH{\,{\buildrel d \over \rightarrow}\,} 
\def\indist{\,{\indistHIGH}\,}

\usepackage[many]{tcolorbox} 
\definecolor{main}{HTML}{3F6EA7}      % primary accent
\definecolor{sub}{HTML}{F2F7FD}       % cool background
\definecolor{sub2}{HTML}{FFF5E8}      % warm background
\definecolor{mainline}{HTML}{D3DFEE}  % soft frame for information boxes
\definecolor{warn}{HTML}{C7771B}      % accent for caution boxes
\definecolor{warnline}{HTML}{EAD9C0}  % soft frame for caution boxes

\tcbset{
    enhanced,
    breakable,
    colback = white,
    boxrule = 0.5pt,
    arc = 2pt,
    left = 9pt,
    right = 9pt,
    top = 7pt,
    bottom = 7pt,
    before skip = 0.3cm,
    after skip = 0.45cm
}                           % setting global options for tcolorbox

\newtcolorbox{boxD}{
    colback = sub,
    colframe = mainline,
    borderline west = {2.2pt}{0pt}{main},
    borderline north = {0.6pt}{0pt}{main!25},
    borderline south = {0.6pt}{0pt}{main!15}
}


\newtcolorbox{boxF}{
    colback = sub2,
    colframe = warnline,
    borderline west = {2.2pt}{0pt}{warn},
    borderline north = {0.6pt}{0pt}{warn!25},
    borderline south = {0.6pt}{0pt}{warn!15}
}

\usepackage{colortbl}


\usepackage{tikz}
\usepackage{verbatim}
\usetikzlibrary{positioning}
\usetikzlibrary{snakes}
\usetikzlibrary{calc}
\usetikzlibrary{arrows}
\usetikzlibrary{decorations.markings}
\usetikzlibrary{shapes.misc}
\usetikzlibrary{matrix,shapes,arrows,fit,tikzmark}




\title{S181M116 Derivatives and Alternative Investments: Lecture 4}
\author{Swapnil Singh}
\date{}

\hypersetup{
  pdftitle={S181M116 Derivatives and Alternative Investments: Lecture 4},
  pdfauthor={Swapnil Singh},
  pdfkeywords={derivatives, options, put-call parity, option bounds, early exercise, binomial trees, delta hedging}
}

\setlength{\headheight}{26pt}

\newcommand{\St}{S_T}
\newcommand{\Szero}{S_0}
\newcommand{\K}{K}
\newcommand{\T}{T}
\newcommand{\dt}{\Delta t}
\newcommand{\erT}{e^{-rT}}
\newcommand{\erdt}{e^{-r\Delta t}}
\newcommand{\PV}{\mathrm{PV}}
\newcommand{\Prob}{\mathbb{P}}
\newcommand{\EQ}{\mathbb{E}^{Q}}
\newcommand{\pstar}{p^*}
\newcommand{\Div}{D}
\newcommand{\Ncontracts}{N}

\begin{document}
\maketitle

\begin{abstract}
This lecture introduces options as asymmetric derivative contracts and develops
the first no-arbitrage restrictions on their prices. We begin with option market
mechanics, option positions, contract specifications, exchange trading, margin,
clearing, and over-the-counter markets. We then study the determinants of option
prices, upper and lower bounds, put-call parity, and early exercise. The final
part introduces binomial trees, risk-neutral valuation, American exercise, delta,
and the link between tree parameters and volatility. Sections 13.10 and 13.11
from the source material are intentionally omitted.
\end{abstract}

\begin{boxD}
\textbf{Reading.} Hull, Chapters 10, 11, and Chapter 13 through Section 13.9.
The practical goal is to read option contracts and option quotes correctly. The
pricing goal is to see how no-arbitrage turns payoffs into restrictions,
relations, and eventually full valuation methods.
\end{boxD}

\section{Roadmap}

Lectures 1--3 studied forwards, futures, and interest-rate instruments. Those
contracts are symmetric commitments: both sides are obligated to transact later.
Options introduce a different structure. The holder has a right but not an
obligation, while the writer has a contingent obligation.

The lecture has three blocks:
\begin{enumerate}
    \item \textbf{Option contracts and markets.} Define calls, puts, long and
    short positions, moneyness, contract specifications, margin, clearing, and
    trading arrangements.
    \item \textbf{Static no-arbitrage restrictions.} Derive bounds and parity
    relations before introducing a full probability model.
    \item \textbf{Binomial valuation.} Price options by building riskless
    portfolios and by using risk-neutral probabilities on one-step and
    multi-step trees.
\end{enumerate}

\begin{boxF}
\textbf{Main lesson.} The option premium is not the expected payoff discounted
at an arbitrary risky discount rate. No-arbitrage prices options by comparing
them with dynamically chosen positions in the underlying asset and the risk-free
asset.
\end{boxF}

\section{Option Basics}

\subsection{Rights versus Obligations}

An option gives its holder a choice. A forward or futures contract locks both
parties into an exchange, but an option allows the holder to walk away when
exercise is unfavorable.

\begin{defN}[Call option]
A call option gives the holder the right to buy the underlying asset at a fixed
strike price $\K$ on or before a fixed maturity date.
\end{defN}

\begin{defN}[Put option]
A put option gives the holder the right to sell the underlying asset at a fixed
strike price $\K$ on or before a fixed maturity date.
\end{defN}

\begin{defN}[Strike and maturity]
The \emph{strike price}, or exercise price, is the price at which the underlying
asset can be bought or sold if the option is exercised. The \emph{maturity} or
expiration date is the last date relevant for exercise.
\end{defN}

The right is valuable because exercise is selective. A call holder exercises only
when buying at $\K$ is better than buying in the market. A put holder exercises
only when selling at $\K$ is better than selling in the market. The writer
receives the premium up front, but must honor the holder's choice.

\subsection{European and American Exercise}

\begin{defN}[European option]
A European option can be exercised only at maturity.
\end{defN}

\begin{defN}[American option]
An American option can be exercised at any time up to and including maturity.
\end{defN}

These names describe exercise style, not geography. European options are easier
to analyze because the exercise date is fixed. American options are usually more
valuable, or at least no less valuable, because they include all European
exercise opportunities plus the possibility of early exercise.

\[
    C \geq c,
    \qquad
    P \geq p,
\]
where $C$ and $P$ denote American call and put prices, and $c$ and $p$ denote
European call and put prices with the same underlying, strike, and maturity.

\subsection{Payoff versus Profit}

The payoff is the value received at expiration before subtracting the initial
premium. The profit subtracts the initial premium, usually ignoring the time
value of money in payoff diagrams.

For a European call:
\[
    \text{payoff}=\max(\St-\K,0).
\]
For a European put:
\[
    \text{payoff}=\max(\K-\St,0).
\]

If the call premium is $c_0$, the call profit at maturity is
\[
    \max(\St-\K,0)-c_0.
\]
If the put premium is $p_0$, the put profit at maturity is
\[
    \max(\K-\St,0)-p_0.
\]

\begin{boxD}
\textbf{Exercise can reduce a loss.} Suppose a call has strike 100 and was
bought for 5. If the stock finishes at 102, the holder exercises. The payoff is
2 and the total profit is $-3$. Exercising does not make the trade profitable,
but it is better than letting the option expire worthless and losing 5.
\end{boxD}

\section{The Four Basic Option Positions}

Every option has a buyer and a writer. The buyer is long the option; the writer
is short the option.

\subsection{Long Call}

A long call benefits from increases in the underlying price. Its payoff is
\[
    \max(\St-\K,0).
\]
The downside is limited to the premium. The upside is potentially unlimited for
a stock because the stock price can in principle rise without bound.

\subsection{Short Call}

A short call is the writer's side of the same contract:
\[
    -\max(\St-\K,0)=\min(\K-\St,0).
\]
The writer receives the premium, but faces losses when the underlying price
rises above the strike. A naked short call is therefore a high-risk position.

\subsection{Long Put}

A long put benefits from decreases in the underlying price:
\[
    \max(\K-\St,0).
\]
The downside is limited to the premium. The payoff is largest when the
underlying price is very low. For a stock, the maximum payoff is approximately
$\K$ because the stock price cannot fall below zero.

\subsection{Short Put}

A short put has payoff
\[
    -\max(\K-\St,0)=\min(\St-\K,0).
\]
The writer receives the premium and loses when the underlying price falls below
the strike. The writer may be forced to buy the asset at $\K$ even when its
market value is much lower.

\begin{boxF}
\textbf{Mnemonic.} Long option positions have rights and pay premiums. Short
option positions receive premiums and carry contingent obligations.
\end{boxF}

\section{Underlying Assets and Contract Design}

\subsection{Stock Options}

Exchange-traded stock options are standardized contracts. In the United States,
one stock option contract typically covers 100 shares. A quoted option price is
usually the price per share, so a quote of 4.20 corresponds to a premium of
\[
    4.20 \times 100 = 420
\]
for one contract.

\subsection{Other Underlyings}

Options can be written on many assets:
\begin{description}
    \item[Exchange-traded products.] Options on ETFs and other listed products
    allow investors to take option exposure to broad equity, bond, commodity, or
    currency baskets.
    \item[Foreign currencies.] Currency options give the right to exchange one
    currency for another at a specified exchange rate.
    \item[Stock indices.] Index options settle in cash because delivering the
    entire index portfolio would be inconvenient.
    \item[Futures contracts.] A futures option gives the right to enter a
    futures position at the strike. Exercise gains depend on the futures price
    relative to the strike.
\end{description}

\subsection{Expiration Dates and Strike Prices}

An exchange must specify which option contracts trade. For stock options this
includes the underlying stock, exercise style, contract size, expiration date,
strike price, and adjustment rules.

Standard U.S. equity options often expire on the third Friday of the expiration
month, while weekly options expire on other Fridays. Longer-dated equity options
are often called LEAPS. Strike spacing is chosen by exchanges and is typically
finer for low-price stocks and wider for high-price stocks.

\subsection{Option Class and Option Series}

\begin{defN}[Option class]
An option class consists of all options of the same type on the same underlying,
such as all IBM calls or all IBM puts.
\end{defN}

\begin{defN}[Option series]
An option series consists of options of a given class with the same strike and
same expiration date, such as October calls with strike 160.
\end{defN}

This distinction matters operationally. A trader may discuss the whole option
class when talking about liquidity in a stock's options, but a trade is in a
specific option series.

\subsection{Moneyness}

Moneyness describes whether immediate exercise would produce a positive payoff.
For stock price $S$ and strike $\K$:
\begin{center}
\begin{tabular}{lll}
\toprule
 & Call & Put \\
\midrule
In the money & $S>\K$ & $S<\K$ \\
At the money & $S=\K$ & $S=\K$ \\
Out of the money & $S<\K$ & $S>\K$ \\
\bottomrule
\end{tabular}
\end{center}

\begin{defN}[Intrinsic value]
The intrinsic value is the payoff from immediate exercise:
\[
    \begin{aligned}
    \text{call intrinsic value} &= \max(S-\K,0),\\
    \text{put intrinsic value}  &= \max(\K-S,0).
    \end{aligned}
\]
\end{defN}

\begin{defN}[Time value]
The time value of an option is its market price minus its intrinsic value.
\end{defN}

An American option cannot sell for less than its intrinsic value because the
holder can exercise immediately. A European option need not satisfy this
immediate-exercise bound because immediate exercise is not allowed.

\subsection{Dividends, Stock Splits, and Contract Adjustments}

Exchange-traded stock options are usually not adjusted for ordinary cash
dividends. The stock price falls on the ex-dividend date, and option prices
reflect expected dividends through valuation.

By contrast, stock splits and stock dividends generally lead to contract
adjustments. The goal is to keep the economic position of the option buyer and
writer unchanged.

\begin{exmp}[Two-for-one stock split]
A call option gives the holder the right to buy 100 shares at 30. After a
two-for-one stock split, the option is adjusted to cover 200 shares at a strike
of 15. The total exercise cost remains
\[
    100 \times 30 = 200 \times 15 = 3000.
\]
\end{exmp}

\begin{exmp}[Stock dividend]
A put option gives the holder the right to sell 100 shares at 15. A 25\% stock
dividend is equivalent to a 5-for-4 split. The adjusted option covers 125 shares
at a strike of 12, keeping total exercise proceeds equal to
\[
    100 \times 15 = 125 \times 12 = 1500.
\]
\end{exmp}

\section{Trading, Margin, Clearing, and Regulation}

\subsection{Market Makers and Bid-Ask Spreads}

Option markets require liquidity across many strikes and maturities. Market
makers quote bid and ask prices and stand ready to trade at those quotes, within
specified size limits. They earn compensation through the bid-ask spread and
manage inventory risk by hedging.

\[
    \text{bid} \leq \text{mid} \leq \text{ask}.
\]

For illiquid option series, the spread can be large relative to the option's
premium. This is especially important for low-priced out-of-the-money options.

\subsection{Offsetting Trades}

Most traders close option positions by entering an offsetting order rather than
by exercising. A trader who bought a call can sell the same call series. A trader
who wrote a put can buy back the same put series. The economic exposure is then
closed through the market.

\subsection{Trading Costs}

Option trades involve more than the model price:
\begin{itemize}
    \item bid-ask spreads;
    \item brokerage commissions and exchange fees;
    \item market impact for large orders;
    \item financing and collateral costs;
    \item tax and accounting treatment, where relevant.
\end{itemize}

Trading costs are small in many textbook derivations, but they matter for
strategy design and for deciding whether an apparent arbitrage can actually be
implemented.

\subsection{Margin for Option Writers}

Buying an option creates no future obligation beyond the premium already paid.
Writing an option creates a possible future obligation, so brokers and clearing
organizations require margin from writers.

For a naked option, the margin protects against adverse movements in the
underlying. For covered or spread positions, margin can be lower because another
position partly offsets the written option's exposure.

\begin{boxF}
\textbf{Economic interpretation.} Margin is not the option's price. It is
collateral against future performance. The writer keeps the option premium but
must post enough collateral to reduce counterparty and clearing risk.
\end{boxF}

\subsection{The Options Clearing Corporation}

For exchange-traded options in the United States, the Options Clearing
Corporation acts as the central clearing counterparty. It records positions,
handles exercises, and guarantees performance by standing between buyers and
writers after trades are cleared.

This clearing structure makes exchange-traded options different from bilateral
over-the-counter options. In an OTC contract, the option holder faces the credit
risk of the writer unless the contract is collateralized or centrally cleared.

\subsection{Exercise Assignment}

When an option holder exercises, the clearing system assigns the exercise to a
writer. The writer usually does not know in advance whether a particular short
option will be assigned, only that assignment can occur if the option is
exercised.

\subsection{Position and Exercise Limits}

Exchanges impose position limits and exercise limits to reduce the chance that a
single trader can dominate a market or create disorderly exercise activity.
Long calls and short puts are typically treated as one side of the market
because both benefit from price increases. Short calls and long puts are treated
as the opposite side.

\subsection{Taxation and Related Instruments}

Tax rules can materially affect option trading incentives. Wash sale rules,
constructive sale rules, and the classification of option gains can matter for
real investors. These rules are jurisdiction-specific and can change, so they
should be treated as institutional background rather than as universal pricing
logic.

Options are also related to warrants, employee stock options, and convertible
bonds:
\begin{description}
    \item[Warrants] are call-like securities issued by a company. Exercise
    usually creates new shares and can dilute existing shareholders.
    \item[Employee stock options] are compensation contracts designed to align
    incentives. They are often at-the-money at grant and subject to vesting
    restrictions.
    \item[Convertible bonds] combine debt with an embedded option to convert
    into equity.
\end{description}

\subsection{Over-the-Counter Options}

OTC options are negotiated bilaterally and can be tailored to a client's exact
needs. They may differ from exchange-traded options in strike, maturity,
underlying, exercise dates, payoff shape, collateral terms, and settlement
rules. Tailoring is valuable, but it introduces greater counterparty credit,
documentation, model, and liquidity risk.

\section{Factors Affecting Option Prices}

\subsection{The Six Main Inputs}

The value of a stock option depends on:
\begin{enumerate}
    \item current stock price $S_0$;
    \item strike price $\K$;
    \item time to expiration $\T$;
    \item stock price volatility $\sigma$;
    \item risk-free interest rate $r$;
    \item dividends expected during the option's life.
\end{enumerate}

These inputs affect calls and puts differently.

\begin{center}
\begin{tabular}{lll}
\toprule
Increase in input & Call value & Put value \\
\midrule
Current stock price $S_0$ & Increases & Decreases \\
Strike price $\K$ & Decreases & Increases \\
Volatility $\sigma$ & Increases & Increases \\
Risk-free rate $r$ & Usually increases & Usually decreases \\
Expected dividends & Decreases & Increases \\
\bottomrule
\end{tabular}
\end{center}

\subsection{Stock Price and Strike Price}

A call is more valuable when the stock price is higher and less valuable when
the strike is higher:
\[
    \max(S_T-K,0)
\]
increases with $S_T$ and decreases with $K$.

A put has the opposite directional exposure:
\[
    \max(K-S_T,0)
\]
increases with $K$ and decreases with $S_T$.

\subsection{Time to Expiration}

For American options, more time cannot reduce value because the longer-dated
option contains every exercise opportunity available to the shorter-dated
option. Thus, for otherwise identical American options, more time generally
means a weakly higher price.

For European options, the effect of time is subtler. A longer-maturity European
option does not permit exercise at the shorter maturity. Large expected
dividends or interest-rate effects can make a shorter European option more
valuable in special cases.

\subsection{Volatility}

Volatility increases both call and put values. An option has limited downside
for the holder but retains upside from favorable movements. Higher volatility
therefore increases the chance that the option finishes deep in the money.

This does not mean volatility makes the underlying asset itself more valuable.
For a stock holder, upside and downside movements both matter. For an option
holder, the payoff truncates the bad side.

\subsection{Risk-Free Interest Rate}

Holding the stock price fixed, a higher risk-free rate tends to increase call
values and decrease put values. Intuitively:
\begin{itemize}
    \item a call delays payment of the strike, so a higher rate makes that
    delayed payment less costly in present-value terms;
    \item a put delivers the strike in the future, so a higher rate reduces the
    present value of that future receipt.
\end{itemize}

\subsection{Dividends}

Expected dividends reduce the stock price when the stock goes ex-dividend. This
hurts calls and benefits puts, all else equal. Dividend timing is also central
to early exercise decisions for American calls.

\section{Assumptions and Notation for No-Arbitrage Results}

We use the following notation:
\[
\begin{aligned}
    S_0 & = \text{current stock price},\\
    K & = \text{strike price},\\
    T & = \text{time to maturity},\\
    r & = \text{continuously compounded risk-free rate},\\
    c & = \text{European call price},\\
    p & = \text{European put price},\\
    C & = \text{American call price},\\
    P & = \text{American put price}.
\end{aligned}
\]

The standard static no-arbitrage arguments assume that large market
participants can borrow and lend at the risk-free rate, short sell where needed,
trade without material transaction costs, and exploit arbitrage opportunities
quickly.

\begin{boxF}
\textbf{Role of assumptions.} These assumptions are not claims that markets are
frictionless in every detail. They identify benchmark restrictions. If a price
violates them by more than trading costs and constraints, arbitrage pressure
should push it back.
\end{boxF}

\section{Upper and Lower Bounds}

\subsection{Upper Bounds}

A call option cannot be worth more than the underlying stock:
\[
    c \leq S_0,
    \qquad
    C \leq S_0.
\]
The call is a right to buy the stock, not ownership of more than the stock.

A European or American put cannot be worth more than the strike:
\[
    p \leq K,
    \qquad
    P \leq K.
\]
For a European put, a sharper bound is
\[
    p \leq K e^{-rT},
\]
because the maximum payoff $K$ occurs at maturity and must be discounted.

\subsection{Lower Bound for a European Call}

For a non-dividend-paying stock, a European call satisfies
\[
    c \geq S_0 - K e^{-rT}.
\]
It also cannot have negative value, so
\[
    c \geq \max(S_0 - K e^{-rT},0).
\]

\begin{proof}
Compare two portfolios:
\begin{description}
    \item[Portfolio A:] one European call plus cash $K e^{-rT}$;
    \item[Portfolio B:] one share of stock.
\end{description}
At maturity, Portfolio A is worth
\[
    \max(S_T-K,0)+K.
\]
If $S_T>K$, this is $S_T$. If $S_T\leq K$, this is $K$, which is at least
$S_T$. Portfolio A is therefore always worth at least as much as Portfolio B at
maturity. It must be worth at least as much today:
\[
    c + K e^{-rT} \geq S_0.
\]
\end{proof}

\begin{exmp}[Call lower bound]
Suppose $S_0=20$, $K=18$, $r=10\%$, and $T=1$. Then
\[
    c \geq 20-18e^{-0.10}
    \approx 20-16.29
    =3.71.
\]
If the call traded for less than 3.71, buying the call and lending the present
value of the strike would dominate owning the stock.
\end{exmp}

\subsection{Lower Bound for a European Put}

For a non-dividend-paying stock, a European put satisfies
\[
    p \geq K e^{-rT}-S_0.
\]
Combining with non-negativity:
\[
    p \geq \max(K e^{-rT}-S_0,0).
\]

\begin{proof}
Compare:
\begin{description}
    \item[Portfolio A:] one European put plus one share;
    \item[Portfolio B:] cash $K e^{-rT}$.
\end{description}
At maturity, Portfolio A is
\[
    \max(K-S_T,0)+S_T.
\]
If $S_T<K$, this equals $K$. If $S_T\geq K$, this equals $S_T$, which is at
least $K$. Portfolio A dominates a risk-free payoff of $K$, so
\[
    p+S_0\geq K e^{-rT}.
\]
\end{proof}

\begin{exmp}[Put lower bound]
Suppose $S_0=37$, $K=40$, $r=5\%$, and $T=0.5$. Then
\[
    p \geq 40e^{-0.05(0.5)}-37
    \approx 39.01-37
    =2.01.
\]
\end{exmp}

\section{Put-Call Parity}

\subsection{European Parity for Non-Dividend-Paying Stocks}

Put-call parity is the central static-arbitrage relation for European options:
\[
    c + K e^{-rT} = p + S_0.
\]

The left side is a fiduciary call: a call plus enough cash to pay the strike at
maturity. The right side is a protective put: a stock plus a put. Both portfolios
pay exactly
\[
    \max(S_T,K)
\]
at maturity.

\begin{center}
\begin{tabular}{llll}
\toprule
State at $T$ & Call + cash & Stock + put & Common payoff \\
\midrule
$S_T>K$ & $S_T-K+K$ & $S_T+0$ & $S_T$ \\
$S_T\leq K$ & $0+K$ & $S_T+K-S_T$ & $K$ \\
\bottomrule
\end{tabular}
\end{center}

Because the terminal payoffs are identical, the current values must be identical
in an arbitrage-free market.

\subsection{Using Parity to Detect Mispricing}

Parity can be rewritten as
\[
    c-p = S_0 - K e^{-rT}.
\]
The difference between the call and put price equals the forward-like value of
buying the stock and financing the strike.

\begin{exmp}[Parity-implied put price]
Suppose $S_0=31$, $K=30$, $r=6\%$, $T=0.25$, and the European call costs 3.
Then parity gives
\[
    p = c + K e^{-rT} - S_0
      = 3 + 30e^{-0.06(0.25)} - 31
      \approx 1.55.
\]
\end{exmp}

\subsection{Arbitrage Logic When Parity Fails}

If
\[
    c+K e^{-rT} > p+S_0,
\]
the fiduciary call is too expensive relative to the protective put. An arbitrage
trader sells the expensive side and buys the cheap side:
\[
    \text{short call, borrow } K e^{-rT}, \text{ buy put, buy stock}.
\]

If
\[
    c+K e^{-rT} < p+S_0,
\]
the protective put is too expensive relative to the fiduciary call. The trader
buys the fiduciary call and sells the protective put:
\[
    \text{buy call, lend } K e^{-rT}, \text{ short put, short stock}.
\]

\subsection{American Options and Parity Inequalities}

American options allow early exercise, so exact European parity does not
generally hold. For a non-dividend-paying stock, the following inequality is the
standard restriction:
\[
    S_0-K \leq C-P \leq S_0-K e^{-rT}.
\]

The upper bound is related to the European parity relation and the fact that an
American put is at least as valuable as a European put. The lower bound reflects
the immediate-exercise value embedded in American contracts.

\section{Early Exercise}

\subsection{American Calls on Non-Dividend-Paying Stocks}

It is never optimal to exercise an American call early when the stock pays no
dividends. The holder gives up two valuable things by exercising:
\begin{enumerate}
    \item the time value of the option;
    \item the interest that could be earned by delaying payment of the strike.
\end{enumerate}

The call holder can participate in upside without paying the strike until
exercise. Exercising early converts the option into stock and unnecessarily pays
$K$ before it is required.

\[
    C = c
\]
for a non-dividend-paying stock when the call is otherwise identical.

\subsection{American Puts on Non-Dividend-Paying Stocks}

Early exercise of an American put can be optimal. A deep-in-the-money put pays
approximately $K-S$ when exercised. If $S$ is very low, exercising allows the
holder to receive $K$ early and earn interest on it.

This early-exercise possibility makes American puts more valuable than European
puts in some cases:
\[
    P \geq p.
\]

The American put lower bound is its intrinsic value:
\[
    P \geq \max(K-S_0,0).
\]

\subsection{Effect of Dividends on Early Exercise}

Dividends can make early exercise of American calls optimal. A call holder does
not receive dividends unless the option is exercised and converted into stock.
Just before a large ex-dividend date, a deep-in-the-money American call may be
exercised to capture the dividend.

For puts, dividends make the underlying stock less valuable after the
ex-dividend date, which can reduce the incentive to exercise immediately before
the dividend. The exact early-exercise decision depends on interest rates,
dividends, moneyness, and remaining time value.

\section{Dividends and Static Bounds}

\subsection{Known Cash Dividends}

Let $D$ be the present value of known cash dividends paid during the option's
life. A stock plus its dividends is economically different from a stock without
dividends. The no-arbitrage bounds adjust by subtracting $D$ from the stock
component.

For European calls:
\[
    c \geq \max(S_0-D-K e^{-rT},0).
\]

For European puts:
\[
    p \geq \max(K e^{-rT}+D-S_0,0).
\]

\subsection{Put-Call Parity with Known Dividends}

With known dividends, European put-call parity becomes
\[
    c + D + K e^{-rT} = p + S_0,
\]
or equivalently
\[
    c-p = S_0-D-K e^{-rT}.
\]

\begin{boxF}
\textbf{Interpretation.} The dividend-adjusted stock price $S_0-D$ is the part
of the stock value that remains attached to the share price at option maturity.
The dividends paid before maturity are not received by the option holder unless
the option is exercised.
\end{boxF}

\section{One-Step Binomial Model}

\subsection{Tree Setup}

The binomial model assumes that over a short period the stock price has only two
possible movements. Starting from $S_0$, the stock moves to
\[
    S_u=S_0u
    \qquad \text{or} \qquad
    S_d=S_0d,
\]
where
\[
    u>1,
    \qquad
    d<1.
\]

Let the option value after an up move be $f_u$ and after a down move be $f_d$.
The goal is to compute the current option value $f$.

\subsection{Riskless Portfolio Argument}

Construct a portfolio with:
\[
    \Delta \text{ shares of stock}
    \quad \text{and} \quad
    -1 \text{ option}.
\]

Its value after an up move is
\[
    S_0u\Delta - f_u.
\]
Its value after a down move is
\[
    S_0d\Delta - f_d.
\]

Choose $\Delta$ so the two values are equal:
\[
    S_0u\Delta - f_u
    =
    S_0d\Delta - f_d.
\]
Solving gives
\[
    \Delta
    =
    \frac{f_u-f_d}{S_0u-S_0d}.
\]

\begin{defN}[Delta in a one-step tree]
Delta is the number of shares needed per option shorted to make the
stock-option portfolio locally riskless:
\[
    \Delta
    =
    \frac{\text{change in option value}}{\text{change in stock value}}.
\]
\end{defN}

Because the portfolio is riskless, it must earn the risk-free rate. Discounting
the certain terminal value back to today determines $f$.

\subsection{Risk-Neutral Probability}

The same value can be written in risk-neutral form:
\[
    f = e^{-rT}\left[p f_u+(1-p)f_d\right],
\]
where
\[
    p=\frac{e^{rT}-d}{u-d}.
\]

The parameter $p$ is not generally the real-world probability of an up move. It
is the probability that makes the expected stock return equal to the risk-free
rate:
\[
    pS_0u+(1-p)S_0d = S_0e^{rT}.
\]

\begin{boxD}
\textbf{Risk-neutral valuation.} To price the option, replace the stock's
expected return by the risk-free rate, compute the expected option payoff under
that adjusted probability, and discount at the risk-free rate.
\end{boxD}

\subsection{Numerical One-Step Example}

Suppose
\[
    S_0=20,\quad u=1.10,\quad d=0.90,\quad K=21,\quad r=12\%,
    \quad T=0.25.
\]
The possible stock prices are
\[
    S_u=22,
    \qquad
    S_d=18.
\]
For a European call,
\[
    f_u=\max(22-21,0)=1,
    \qquad
    f_d=\max(18-21,0)=0.
\]
The risk-neutral probability is
\[
    p=
    \frac{e^{0.12(0.25)}-0.90}{1.10-0.90}
    \approx
    \frac{1.03045-0.90}{0.20}
    =0.6523.
\]
The option value is
\[
    f=e^{-0.12(0.25)}[0.6523(1)+0.3477(0)]
    \approx 0.633.
\]

The hedge ratio is
\[
    \Delta=\frac{1-0}{22-18}=0.25.
\]
Thus a trader who is short one call hedges the one-step risk by buying 0.25
shares.

\section{Why the Real-World Expected Return Drops Out}

The option is risky, so it is tempting to discount its expected payoff using a
risky expected return. The binomial argument shows why this is not the right
primitive. The option is priced relative to the stock. Once the stock price
$S_0$ is observed, the stock already embeds the market's required compensation
for risk.

The no-arbitrage relation between the option and the stock does not require the
stock's expected real-world return $\mu$. Different investors may disagree about
$\mu$, but if they agree on $S_0$, $u$, $d$, $r$, and the payoff, the
arbitrage-free option price is pinned down.

\begin{boxF}
\textbf{Key distinction.} Risk-neutral valuation does not say investors are
risk-neutral in reality. It says that for pricing derivatives relative to traded
underlyings, we may use risk-neutral probabilities because the risk adjustment is
already reflected in traded asset prices.
\end{boxF}

\section{Two-Step Binomial Trees}

\subsection{Backward Induction}

In a two-step tree, the stock can move up or down in each step. Starting from
$S_0$, the terminal stock prices are
\[
    S_0u^2,
    \qquad
    S_0ud,
    \qquad
    S_0d^2.
\]
If $ud=du$, the middle node is recombining: one up then one down gives the same
stock price as one down then one up.

The valuation procedure is:
\begin{enumerate}
    \item compute terminal option payoffs;
    \item value the option one step earlier using the one-step formula;
    \item continue backward until reaching the initial node.
\end{enumerate}

For time step $\dt$, the one-step recursion is
\[
    f=e^{-r\dt}\left[p f_u+(1-p)f_d\right],
\]
where
\[
    p=\frac{e^{r\dt}-d}{u-d}.
\]

\subsection{Two-Step Formula for a European Option}

Let the terminal option payoffs be $f_{uu}$, $f_{ud}$, and $f_{dd}$. Then
\[
\begin{aligned}
    f_u &= e^{-r\dt}\left[p f_{uu}+(1-p)f_{ud}\right],\\
    f_d &= e^{-r\dt}\left[p f_{ud}+(1-p)f_{dd}\right],\\
    f   &= e^{-r\dt}\left[p f_u+(1-p)f_d\right].
\end{aligned}
\]
Substituting gives
\[
    f=e^{-2r\dt}
    \left[
        p^2 f_{uu}
        +2p(1-p)f_{ud}
        +(1-p)^2 f_{dd}
    \right].
\]

\subsection{European Put Example}

Consider a two-year European put with
\[
    S_0=50,\quad K=52,\quad u=1.2,\quad d=0.8,\quad r=5\%,
    \quad \dt=1.
\]
The risk-neutral probability is
\[
    p=\frac{e^{0.05}-0.8}{1.2-0.8}
    \approx 0.6282.
\]
The terminal stock prices are
\[
    72,\quad 48,\quad 32.
\]
The terminal put payoffs are
\[
    f_{uu}=0,\quad f_{ud}=4,\quad f_{dd}=20.
\]
Thus
\[
\begin{aligned}
    f
    &=
    e^{-0.10}
    \left[
        (0.6282)^2(0)
        +2(0.6282)(0.3718)(4)
        +(0.3718)^2(20)
    \right] \\
    &\approx
    4.19.
\end{aligned}
\]

\section{American Options in Trees}

\subsection{Exercise Check at Each Node}

For American options, backward induction must include an exercise decision at
every node before maturity. At a node with stock price $S$, the value is
\[
    \max\{\text{intrinsic value},\text{continuation value}\}.
\]

For an American put:
\[
    \text{node value}
    =
    \max\left[
        K-S,\,
        e^{-r\dt}\left(p f_u+(1-p)f_d\right)
    \right].
\]

For an American call:
\[
    \text{node value}
    =
    \max\left[
        S-K,\,
        e^{-r\dt}\left(p f_u+(1-p)f_d\right)
    \right].
\]

\subsection{Early Exercise Example Logic}

Return to the put with $S_0=50$, $K=52$, $u=1.2$, $d=0.8$, $r=5\%$, and
$\dt=1$. At the lower node after one down move, the stock price is 40. Immediate
exercise gives
\[
    K-S=52-40=12.
\]
The continuation value is the discounted risk-neutral expectation of next
period's values. If continuation is less than 12, the American put is exercised
at that node.

\begin{boxD}
\textbf{Backward induction with American exercise.} First compute continuation
value as if the option were European from that node onward. Then compare it with
immediate exercise. The node value is the larger of the two.
\end{boxD}

\section{Delta Hedging}

\subsection{Delta as Local Sensitivity}

Delta measures the sensitivity of the option price to the underlying price:
\[
    \Delta
    =
    \frac{\partial f}{\partial S}
\]
in continuous models, and
\[
    \Delta
    =
    \frac{f_u-f_d}{S_u-S_d}
\]
in a one-step binomial tree.

Call deltas are positive. Put deltas are negative. A deeply in-the-money call
has delta close to 1, while a deeply out-of-the-money call has delta close to 0.
A deeply in-the-money put has delta close to $-1$, while a deeply
out-of-the-money put has delta close to 0.

\subsection{Delta Changes Over Time}

In a multi-step tree, delta is not constant. It differs by node because the
option's moneyness and remaining time differ by node. A hedge that is riskless
over one small interval must be rebalanced as the stock price and time to
maturity change.

\begin{boxF}
\textbf{Dynamic hedging.} Option hedging is dynamic. The portfolio that is
riskless over one short interval is not usually riskless over the next interval
unless the stock position is adjusted.
\end{boxF}

\section{Matching Binomial Trees to Volatility}

\subsection{Choosing $u$ and $d$}

A binomial tree needs three parameters per time step:
\[
    u,\quad d,\quad p.
\]
Once $u$ and $d$ are chosen, no-arbitrage fixes $p$:
\[
    p=\frac{e^{r\dt}-d}{u-d}.
\]

The up and down multipliers are chosen to match volatility. If annual volatility
is $\sigma$ and the time step is $\dt$, the standard binomial choice is
\[
    u=e^{\sigma\sqrt{\dt}},
    \qquad
    d=e^{-\sigma\sqrt{\dt}}=\frac{1}{u}.
\]

Then
\[
    p=\frac{e^{r\dt}-d}{u-d}.
\]

\subsection{Why Volatility Is Preserved Under Risk-Neutral Valuation}

Moving from the real-world measure to the risk-neutral measure changes expected
returns. The real-world expected stock growth rate $\mu$ is replaced by the
risk-free rate $r$. Volatility is not changed. The tree should therefore match
the asset's volatility under both the real-world and risk-neutral measures, but
use a risk-neutral probability when pricing.

\begin{boxD}
\textbf{Connection to measure changes.} In more advanced notation, the
real-world probability measure is often called $P$, while the risk-neutral
measure is called $Q$. The change from $P$ to $Q$ changes drifts, not the local
volatility parameter used to build the tree.
\end{boxD}

\subsection{Numerical Tree Parameters}

Suppose a two-year option is valued with two time steps, so $\dt=1$. Let
\[
    \sigma=30\%,
    \qquad
    r=5\%.
\]
Then
\[
    u=e^{0.30}=1.3499,
    \qquad
    d=e^{-0.30}=0.7408,
\]
and
\[
    p=
    \frac{e^{0.05}-0.7408}{1.3499-0.7408}
    \approx
    0.5097.
\]

These parameters define the tree. For an American option, each node then uses
the continuation-versus-exercise comparison.

\section{Increasing the Number of Steps}

\subsection{More Steps Means a Finer Approximation}

The same tree formulas apply for any number of steps $n$. If the option life is
$T$, then
\[
    \dt=\frac{T}{n}.
\]
The parameters are
\[
    u=e^{\sigma\sqrt{\dt}},
    \qquad
    d=e^{-\sigma\sqrt{\dt}},
    \qquad
    p=\frac{e^{r\dt}-d}{u-d}.
\]

As $n$ increases, the tree has a finer time grid and a richer set of possible
terminal stock prices. For European options on non-dividend-paying stocks, the
binomial price converges to the Black-Scholes-Merton price as the number of
steps becomes large.

\subsection{Five-Step Illustration}

Suppose $T=2$, $n=5$, $r=5\%$, and $\sigma=30\%$. Then
\[
    \dt=\frac{2}{5}=0.4.
\]
The tree parameters are
\[
    u=e^{0.30\sqrt{0.4}}\approx 1.2089,
    \qquad
    d=\frac{1}{u}\approx 0.8272,
\]
\[
    e^{r\dt}=e^{0.05(0.4)}\approx 1.0202,
    \qquad
    p=\frac{1.0202-0.8272}{1.2089-0.8272}\approx 0.5056.
\]

The pricing algorithm does not change. Compute terminal payoffs, work backward
one step at a time, and for American options compare continuation value with
exercise value at each node.

\section{What Is Omitted}

The source material includes Section 13.10 on DerivaGem software and Section
13.11 on options on other assets. These notes intentionally exclude those two
sections. The binomial-tree discussion here is limited to stock options through
the increasing-number-of-steps material in Section 13.9.

\section{Worked Examples}

\begin{exmp}[Payoff and profit of a long call]
A call has strike 100 and premium 6. At maturity the stock price is 114. The
payoff is
\[
    \max(114-100,0)=14.
\]
The profit per share is
\[
    14-6=8.
\]
For one standard contract on 100 shares, the profit is 800.
\end{exmp}

\begin{exmp}[Payoff and profit of a long put]
A put has strike 70 and premium 4. At maturity the stock price is 58. The payoff
is
\[
    \max(70-58,0)=12.
\]
The profit per share is
\[
    12-4=8.
\]
\end{exmp}

\begin{exmp}[Intrinsic and time value]
A call has strike 50, the stock price is 56, and the call price is 8.50. The
intrinsic value is
\[
    \max(56-50,0)=6.
\]
The time value is
\[
    8.50-6=2.50.
\]
\end{exmp}

\begin{exmp}[European parity arbitrage check]
Suppose $S_0=40$, $K=40$, $r=5\%$, $T=1$, $c=4.00$, and $p=2.00$. Parity
requires
\[
    c+Ke^{-rT}
    =
    p+S_0.
\]
The left side is
\[
    4+40e^{-0.05}=42.05.
\]
The right side is
\[
    2+40=42.00.
\]
The apparent mispricing is 0.05 per share before transaction costs. In practice,
this would likely be too small to exploit after bid-ask spreads and fees.
\end{exmp}

\begin{exmp}[American put node decision]
At a node in an American put tree, $S=38$, $K=45$, $r=4\%$, $\dt=0.5$, $p=0.55$,
$f_u=5$, and $f_d=12$. The continuation value is
\[
    e^{-0.04(0.5)}[0.55(5)+0.45(12)]
    =
    e^{-0.02}(8.15)
    \approx 7.99.
\]
Immediate exercise gives
\[
    K-S=45-38=7.
\]
The node value is 7.99, so the put is not exercised at this node.
\end{exmp}

\begin{exmp}[Delta in a tree]
At a node, the stock can move to 66 or 54. A call will be worth 9 after the up
move and 1 after the down move. Its delta is
\[
    \Delta=\frac{9-1}{66-54}=\frac{8}{12}=0.6667.
\]
A trader short 100 calls would buy approximately 66.67 shares to create the
one-step delta hedge.
\end{exmp}

\section{Checklist for Students}

After this lecture you should be able to:
\begin{itemize}
    \item distinguish option rights from forward and futures obligations;
    \item define calls, puts, strike price, expiration, European exercise, and
    American exercise;
    \item compute payoffs and profits for long and short calls and puts;
    \item explain why option writers face margin requirements;
    \item distinguish option class, option series, intrinsic value, time value,
    and moneyness;
    \item describe how stock splits and stock dividends affect option contract
    terms;
    \item explain the roles of exchanges, market makers, brokers, clearing
    members, and the Options Clearing Corporation;
    \item list the main factors affecting stock option prices and sign their
    effects on calls and puts;
    \item derive and use upper and lower bounds for European and American
    options;
    \item derive put-call parity for European options on non-dividend-paying
    stocks;
    \item adapt parity and bounds for known dividends;
    \item explain why an American call on a non-dividend-paying stock should not
    be exercised early;
    \item explain why an American put may be exercised early;
    \item set up a one-step binomial tree and compute the risk-neutral
    probability;
    \item value a European option by risk-neutral expectation on a tree;
    \item construct the one-step delta hedge;
    \item value a two-step European option by backward induction;
    \item value an American option by comparing continuation and exercise values
    at each node;
    \item choose $u$, $d$, and $p$ from volatility, the risk-free rate, and the
    time step;
    \item explain why the real-world expected stock return does not enter the
    no-arbitrage option price;
    \item explain why binomial prices converge as the number of steps increases.
\end{itemize}

\section{Practice Questions}

\begin{enumerate}
    \item A call has strike 60 and premium 4. Compute the payoff and profit when
    the terminal stock price is 52, 60, and 71.
    \item A put has strike 45 and premium 3.50. Compute the payoff and profit
    when the terminal stock price is 35, 45, and 52.
    \item Draw the payoff diagrams for a long call, short call, long put, and
    short put. Label the strike and the break-even points after including
    premiums.
    \item One equity option contract covers 100 shares. If a call is quoted at
    2.75, what premium is paid for 12 contracts?
    \item A stock is at 84. Classify a call with strike 80, a call with strike
    90, a put with strike 80, and a put with strike 90 as in, at, or out of the
    money.
    \item A put with strike 100 is priced at 13 when the stock is 91. Compute
    intrinsic value and time value.
    \item A company has a 3-for-2 stock split. How should the terms of a call on
    100 shares with strike 48 be adjusted?
    \item Explain why a naked short call requires margin but a purchased call
    does not.
    \item List the six main factors affecting stock option prices. For each,
    state whether an increase usually raises or lowers a call and a put.
    \item A non-dividend-paying stock is 72, the strike is 70, $r=4\%$, and
    $T=0.5$. Compute the lower bound for the European call.
    \item A non-dividend-paying stock is 38, the strike is 40, $r=3\%$, and
    $T=1$. Compute the lower bound for the European put.
    \item A European call costs 6.20. The stock is 50, the strike is 52,
    $r=5\%$, and $T=0.75$. Use put-call parity to compute the corresponding
    European put price.
    \item Suppose $c+Ke^{-rT}>p+S_0$. Describe the arbitrage trade and verify
    that the terminal payoff is nonnegative.
    \item Why is it generally suboptimal to exercise an American call early
    when the stock pays no dividends?
    \item Why can it be optimal to exercise a deep-in-the-money American put
    early?
    \item A stock is 30 and can move to 36 or 24 in one year. A one-year call
    has strike 32 and $r=5\%$. Compute $p$ and the call price.
    \item In the previous question, compute the hedge ratio $\Delta$.
    \item A two-step tree has $S_0=100$, $u=1.1$, $d=0.9$, $r=4\%$, and
    $\dt=0.5$. Compute the risk-neutral probability.
    \item Use the tree in the previous question to value a European put with
    strike 100.
    \item At an American put node, $S=42$, $K=50$, $r=6\%$, $\dt=0.25$,
    $p=0.52$, $f_u=5$, and $f_d=11$. Should the option be exercised?
    \item Compute $u$, $d$, and $p$ for a stock option tree when
    $\sigma=25\%$, $r=4\%$, and $\dt=1/12$.
    \item Explain in words why risk-neutral probabilities are not generally the
    same as real-world probabilities.
    \item If the number of tree steps is increased from 10 to 500 for a fixed
    option maturity, what happens to $\dt$? Why does this improve the
    approximation?
    \item State clearly which parts of the source material are omitted from
    these notes and why.
\end{enumerate}

\end{document}
